\documentclass[10pt]{article}

\usepackage[utf8]{inputenc}

\usepackage[T1]{fontenc}

\usepackage{amsmath}

\usepackage{amsfonts}

\usepackage{amssymb}

\usepackage[version=4]{mhchem}

\usepackage{stmaryrd}

\usepackage{graphicx}

\usepackage[export]{adjustbox}

\graphicspath{ {./images/} }

\usepackage{mathrsfs}



\begin{document}

между определяемыми из опыта величинами - дисперсионные соотношения, Померанчука теорему и др.



В случае упругого рассеяния бесспиновых частиц решение Шредингера уравнения для волновой функции $\psi(\boldsymbol{r})$ при $r \rightarrow \infty$ имеет вид:



\[

\left.\psi(\boldsymbol{r})\right|_{r \rightarrow \infty} \sim \exp (i \boldsymbol{r} \boldsymbol{r})+f(\vartheta) r^{-1} \exp (i \boldsymbol{r} \boldsymbol{r}) .

\]



Здесь $r$ - расстояние между частицами, $\boldsymbol{k}=\boldsymbol{p} / \hbar-$ волновой вектор, $\boldsymbol{p}$ - импульс в системе центра инерции сталкивающихся частиц, $\vartheta$ - угол рассеяния, $f(\vartheta)$ - амплитуда рассеяния - функция, зависящая от угла рассеяния и энергии сталкивающихся частиц. Первый член в этом выражении описывает падающие частицы, второй - рассеянные. Дифференциальное сечение рассеяния определяется как отношение числа частиц, рассеянных за единицу времени в элемент телесного угла $d \Omega$, к плотности потока падающих частиц. Сечение рассеяния на угол $\vartheta$ (в системе центра инерции) в единичный телесный угол равно:



\[

\frac{d \sigma}{d \Omega}=|f(\vartheta)|^{2}

\]



Амплитуду рассеяния обычно разлагают в ряд по п а $\mathrm{p}$ ц и аль ным волнам - состояниям с определенным орбитальным моментом $l$ :



\[

f(\vartheta)=\frac{1}{2 i k} \sum_{l=0}^{\infty}(2 l+1)\left(S_{l}-1\right) P_{l}(\cos \vartheta) .

\]



Здесь $P_{l}(\cos \vartheta)$ - многочлены Лежандра, $S_{l}$ - комплексные функщии энергии, зависящие от характера взаимодействия и являющиеся элементами $S$-матрицы (в представлении, в к-ром диагональны энергия, угловой момент и его проекция). Если число падающих на центр частиц с орбитальным моментом $l$ равно числу идущих от центра частиц с тем же моментом (упругое рассеяние), то $\left|S_{/}\right|=1$. В общем случае $\left|S_{l}\right| \leqslant 1$. Эти условия - следствие условия унитарности $S$-матрицы. Если возможно только упругое рассеяние, то $S_{l}=\exp \left(2 i \delta_{l}\right)$ и рассеяние в состоянии с данным $l$ характеризуется только одним вещественным параметром $\delta_{l}-\phi$ а зой рассея н и я. Если $\delta_{l}=0$ при нек-ром $l$, то рассеяние в состоянии с орбитальным моментом $l$ отсутствует.



Полное сечение упругого рассеяния равно:



\[

\begin{gathered}

\sigma^{\mathrm{y} \Pi \mathrm{p}}=\sum_{l=0}^{\infty} \sigma_{l}^{\mathrm{y} п \mathrm{p}}, \\

\sigma_{l}^{\mathrm{y} \Pi \mathrm{p}}=\pi \lambda^{2}(2 l+1)\left|S_{l}-1\right|^{2},

\end{gathered}

\]



где $\sigma_{l}^{\text {упр }}$ - парциальное сечение упругого рассеяния частиц с орбитальным моментом $l, \lambda=1 / k$ - длина волны де Брой-



\includegraphics[max width=\textwidth, center]{2023_07_08_e7d342f458619e694008g-1}

и равно:



\[

\left(\sigma_{l}^{\mathrm{ynp}}\right)_{\text {макс }}=4 \pi \lambda^{2}(2 l+1)

\]



при этом $\delta_{l}=\pi / 2$ (резонанс в рассеянии). Таким образом, при резонансе сечение процесса определяется де-бройлевской длиной волны $\lambda$ и для медленных частиц, для к-рых $\lambda \gg R_{0}$, где $R_{0}-$ радиус действия сил, намного превосходит величину $\pi R_{0}^{2}$ (классич. сечение рассеяния). Это явление (необъяснимое с точки зрения классич. теории рассеяния) обусловлено волновой природой микрочастиц.



Другим проявлением волновой природы микрочастиц является дифракционное рассеяние - упругое рассеяние быстрых частиц на малые углы $\vartheta \sim \lambda / R_{0}$ (при $\lambda \ll R_{0}$ ), обусловленное отклонением де-бройлевских волн налетающих частиц в область геометрич. тени, возникающей за рассеивающей частицей. Таким образом, дифракционное рассеяние аналогично явлению дифракции света.



Зависимость сечения рассеяния от энергии вблизи резонанса определяется формулой Брейта - Вигнера:



\[

\sigma_{l}=4 \pi \lambda^{2}(2 l+1) \frac{(\Gamma / 2)^{2}}{\left(\varepsilon-\varepsilon_{0}\right)+(\Gamma / 2)^{2}},

\]



где $\mathscr{E}_{0}-$ энергия, при к-рой сечение достигает максимума (положение резонанса), а $\Gamma-$ ширина резонанса. При $\mathscr{E}=\mathscr{E}_{0}+\Gamma / 2$ сечение $\sigma_{l}=\sigma_{l}^{\text {макс }} / 2$.



Полное сечение всех неупругих процессов



\[

\begin{gathered}

\sigma^{\text {неупр }}=\sum_{l=0}^{\infty} \sigma_{l}^{\text {неупр }}, \\

\sigma_{l}^{\text {неупр }}=\pi \lambda^{2}(2 l+1)\left(1-\left|S_{l}\right|^{2}\right) .

\end{gathered}

\]



Условие унитарности ограничивает величину парциального сечения для неупругих процессов:



\[

\sigma_{l}^{\text {неупр }} \leqslant \pi \lambda^{2}(2 l+1)

\]



Для короткодействующих потенциалов взаимодействия основную роль играют фазы рассеяния с $l \leqslant R_{0} / \lambda$, где $R_{0}-$ радиус действия сил; величина $l \lambda$ определяет минимальное расстояние, на к-рое может приблизиться к центру сил свободная частица с моментом $l$ (прицельный параметр в квантовой теории). При $R_{0} / \lambda \ll 1$ (малые энергии) следует учитывать только парциальную волну с $l=0$ ( $S$-волну). Амплитуда рассеяния в этом случае



\[

f \approx \frac{1}{2 i k}\left(\exp \left(2 i \delta_{0}\right)-1\right)=\frac{1}{k \operatorname{ctg} \sigma_{0}-i k}

\]



и сечение рассеяния не зависит от $\vartheta$ (рассеяние сферически симметрично). При малых энергиях



\[

k \operatorname{ctg} \sigma_{0} \approx-\frac{1}{a}+\frac{1}{2} r_{0} k^{2} .

\]



Параметры $a$ и $r_{0}$ называются соответственно д л и н о й рассеянияи эффективным радиусом рассеяния. Их находят из опыта, и они являются важными характеристиками сил, действующих между частицами. Длина рассеяния равна по величине и противоположна по знаку амплитуде рассеяния при $k=0$. Полное сечение рассеяния при $k=0$ равно $\sigma_{0}=4 \pi a^{2}$.



Если у частиц имеется связанное состояние с малой энергией связи, то их рассеяние при $R_{0} / \lambda \ll 1$ носит резонансный характер. Типичный пример - рассеяние нейтронов протонами в состоянии с полным спином $J=1$, в к-ром система нейтрон-протон имеет связанное состояние (дейтрон). В этом случае длина рассеяния $\boldsymbol{a}$ отрицательна, а сечение рассеяния зависит только от энергии связи.



Если параметр $R_{0} / \lambda$ невелик, фазы рассеяния могут быть получены из измеренных на опыте сечений, поляризаций и других величин. Эта процедура называется фазовым а нал и зом. Найденные фазы рассеяния сравниваются с теоретич. предсказаниями и позволяют получить важную информацию о характере взаимодействия.



Информацию о взаимодействии дают измерения поляризационных эффектов. Для упругого рассеяния частиц со спином 0 на частицах со спином 1/2 вместо (2) имеем



\[

\psi_{0}(\boldsymbol{r}) \sim \exp (i \boldsymbol{r}) u_{\sigma}+M_{\sigma \sigma^{\prime}}\left(\boldsymbol{k}^{\prime}, \boldsymbol{k}\right) u_{\sigma^{\prime}} \exp (i \boldsymbol{r}) / r .

\]



Здесь $\boldsymbol{k}=\boldsymbol{p} /|\boldsymbol{p}|$ и $\boldsymbol{k}^{\prime}=\boldsymbol{p}^{\prime} /\left|\boldsymbol{p}^{\prime}\right|\left(\boldsymbol{p}\right.$ и $\boldsymbol{p}^{\prime}-$ начальный и конечный импульсы в системе центра инерции), $u_{\sigma}$ - спинор, описывающий состояние начальных частиц, $M_{\sigma \sigma^{\prime}}\left(\boldsymbol{k}_{1}, \boldsymbol{k}\right)$ $-(2 \times 2)$-матрица, называемая с п и но вой м а т р ицей р а с с е я н и я, $\sigma$ - спиновый индекс (по повторяющемуся индексу $\sigma^{\prime}$ производится суммирование). Из сохранения полного момента и четности (инвариантности относительно вращений и пространственных отражений) следует, что матрица $M$ имеет общий вид



\[

M=a+b \sigma \boldsymbol{n},

\]



где $a$ и $b$ - комплексные функции скаляров $\boldsymbol{k} \boldsymbol{k}^{\prime}$ и $\boldsymbol{k}^{2}, \sigma_{i}$ Паули матрицы, $\boldsymbol{n}=\left[\boldsymbol{k} \boldsymbol{k}^{\prime}\right] /\left|\left[\boldsymbol{k} \boldsymbol{k}^{\prime}\right]\right|-$ единичный вектор нормали к плоскости рассеяния.





\end{document}